The welfare effect of a price change and the income response inferred from
purchases are different objects. A household at origin $n$ spends income $v_n$
on two destinations and an outside good. Let $p_{ni}$ be its purchase price,
including monetary access costs, and $c_{ni}$ its chosen quantity. Prices are
positive. Preferences are fixed and the household optimizes over a fixed feasible
set. Utility is strictly increasing in the outside good, so initial purchases
use all income and initial expenditure equals income. An expenditure
function $e_n(p_n,u_n)$ gives the minimum cost of its initial utility. At a
differentiability point with an attained initial optimum, its price derivative
is the initial compensated purchase quantity.

These conditions give a first-order welfare change in money:
\[
 dW_n=dv_n-\sum_i c_{ni}\,dp_{ni},\qquad
 \frac{dW_n}{v_n}=d\log v_n-\sum_i s_{ni}\,d\log p_{ni},
 \quad s_{ni}=\frac{p_{ni}c_{ni}}{v_n}.
\]
The sum includes the outside good, so initial shares add to one. Initial
purchases value price changes even though the household subsequently substitutes.
This local valuation does not yet determine income from observed spending.

Recovering income requires separating a household's resources from the change
in where it buys. Write spending as $X_{ni}=v_n s_{ni}$. To make that separation,
suppose each destination has a common relative share response $D_i$ across
origins, with the outside response normalized to zero:
\[
 d\log s_{ni}=D_i-\sum_j s_{nj}D_j,\qquad D_0=0.
\]
This restriction allows different initial baskets but rules out origin-specific
destination responses. Substitution then gives
\[
 d\log X_{ni}=\alpha_n+D_i,\qquad
 \alpha_n=d\log v_n-\sum_j s_{nj}D_j.
\]
Thus an estimated origin effect is not income itself. Adding the common
destination responses weighted by the initial basket at that origin recovers income
under this restriction and a compatible normalization.

For example, a change in monetary access costs that affects nearby residents
more than distant residents need not generate a common destination response.
The origin effect correction can then fail, while the household envelope remains
valid using its own purchase prices. Estimation therefore needs evidence for
the common response as well as variation separating the two effects. If an
outside response or a common income component is not observed or identified,
the income level across the city cannot be recovered merely by naming the fitted origin term.
